% Propietario conservado: /Users/ruben/Documents/New project/output/REVISION_ARGUMENTAL_APERTURAS_CIERRES_20260915/VII/delivery/MANUSCRIPT_VII_ES_EN_COMPLETE/source_es/manuscrito/33_corriente_espinorial_y_densidad.tex
% Artículo VII. Incorporación aditiva autorizada, 10-09-2026.
% RESULTADO_RECUPERADO: representación real de Clifford y CAR:
% TEORIA_HOLOGRAFICA_INTEGRAL_MASA_HMT_MD_REV2_2026-08-20/fuente/modulos/
% 18_cuantica_estadistica_y_termodinamica.md, líneas 131--181 y 218--255;
% acción de primer orden: 14_topologia_exoticos_y_gravedad.md, 459--517;
% distinción corriente media/segundo momento y dilución:
% HMT_OBRAS/OBRA_II/chapters/11_cartan_bianchi_cosmologia.tex, 201--230.
% PDF REV2: b56bb38d10b9a4f4d0de47b7da8010fb31a6e8e8684a3bfd4237e37d7a3dddfa.
% FORMALIZACION_NUEVA: acción espinorial afín explícita, contracción
% normalizada mediante 31, observable cuártico CAR y funcional de amplitud.
% CERTIFICADO_NUEVO: cálculo racional de la forma cuadrática en las cuatro
% componentes de una trivector y evaluación de la realización CAR finita.
% Convención explicitada: las matrices g tienen firma (-+++); A_D=-i g^0
% y la acción cinética simétrica lleva hbar/2. La combinación equivalente
% i hbar Gamma^I usa Gamma^I=-i g^I y firma opuesta. No se mezclan ambas.
% Dependencias: cotetrada/conexión de 30d; constantes internas heredadas;
% corriente e inversas de 31. Estado material, celda y normal orden son
% datos de la realización declarada, no constantes universales extraídas.
% La acción afín de este fragmento y la acción de extensión mínima de 30e
% conservan sus problemas variacionales propios. No se identifican sus
% corrientes. No se cambia 50: su dominio de amplitud positiva se concreta.

\section{Courant spinoriel et amplitude torsionnelle de l'état matériel}
\label{vii:vii:sec:corriente-espinorial-densidad}

Le développement précédent et la réalisation spinorielle qui suit articulent
deux problèmes variationnels ayant des antécédents géométriques communs. L'
action d'extension minimale \(S_{\min}=\mathfrak a_m E\), construite dans
la section~\ref{vii:vii:sec:corriente-extension-minima}, incorpore la
dépendance de l'incidence transportée et de son minimum à l'égard de la
connexion ; son élimination torsionnelle conserve cette dépendance complète dans
la section~\ref{vii:vii:sec:eliminacion-torsional-no-lineal}. L'action
spinorielle \(S_D\) de \eqref{vii:vii:eq:accion-dirac-material} couple la
même connexion aux champs de sa représentation et est affine en la
contorsion lorsque la cotétrade et les champs sont fixés. Chaque courant s'obtient
en faisant varier sa propre action. Une préparation minimisante peut déterminer
l'état matériel initial de la réalisation spinorielle ; la substitution
de champs dépendant de la connexion conserve en outre les termes
de la différentielle composée. Cette distinction situe le calcul d'amplitude
suivant dans son action et dans son domaine variationnel spécifiques.

La représentation spinorielle du transport et la réponse algébrique de Cartan--Holst permettent de calculer l'intensité torsionnelle comme une fonctionnelle de l'état matériel. Son évaluation conserve trois données : la représentation des champs, le second moment de leur courant et le volume par rapport auquel la densité est définie. Le coefficient de couplage résulte de l'action et de son inversion ; l'état initial détermine l'amplitude matérielle.

\subsection{Représentation spinorielle et action matérielle affine}
\label{vii:vii:subsec:accion-espinorial-afin}

Nous conservons la cotétrade, l'orientation et les couplages de la
section~\ref{vii:vii:sec:corriente-respuesta-cartan}. La réalisation
spinorielle utilise le module complexe de dimension quatre obtenu en
complexifiant les matrices réelles
\begin{equation}
 \begin{gathered}
 J=\begin{pmatrix}0&-1\\1&0\end{pmatrix},\qquad
 R=\begin{pmatrix}0&1\\1&0\end{pmatrix},\qquad
 Z=\begin{pmatrix}1&0\\0&-1\end{pmatrix},\\
 g^0=J\otimes I_2,\qquad g^1=R\otimes I_2,\qquad
 g^2=Z\otimes R,\qquad g^3=Z\otimes Z.
 \end{gathered}
 \label{vii:vii:eq:clifford-material-explicito}
\end{equation}
La multiplication donne
\(\{g^I,g^J\}=2\eta^{IJ}I_4\), avec
\(\eta=\operatorname{diag}(-1,1,1,1)\). Les matrices se distinguent
du paramètre scalaire \(\gamma\) de l'opérateur de Holst. Nous définissons
\begin{equation}
 g_I=\eta_{IJ}g^J,\quad
 \mathbb B_{IJ}=\tfrac14[g_I,g_J],\quad
 A_D=-i g^0,\quad \bar\psi=\psi^\dagger A_D,\quad
 \Omega=\tfrac12\omega^{IJ}\mathbb B_{IJ}.
 \label{vii:vii:eq:adjunto-conexion-espinorial}
\end{equation}
Le produit de Clifford vérifie
\[
 [\mathbb B_{IJ},g^K]=\delta_J^K g_I-\delta_I^K g_J.
\]
Ainsi, la même connexion métrique agit sur les champs par
\(D\psi=\mathrm d\psi+\Omega\psi\) et
\(D\bar\psi=\mathrm d\bar\psi-\bar\psi\Omega\).
Le transport sur une route s'obtient en intégrant cette connexion dans
sa représentation spinorielle ; il conserve la connexion et la route de la
réalisation géométrique précédente.

Soit \(\eta_I=\iota_{E_I}\operatorname{vol}_e\), de sorte que \(e^J\wedge\eta_I=\delta_I^J\operatorname{vol}_e\). Pour un champ de type dimensionnel \([\psi]=L^{-3/2}\), nous prenons l'action matérielle minimale
\begin{equation}
 S_D=\int_M\left\{
 \frac{\hbar}{2}
 \bigl[\bar\psi g^I D\psi-(D\bar\psi)g^I\psi\bigr]\wedge\eta_I
 -mc\,\bar\psi\psi\operatorname{vol}_e\right\}.
 \label{vii:vii:eq:accion-dirac-material}
\end{equation}
Ici, \(\hbar=\hbar_{\rm int}\), \(c=c_{\rm int}\) et la masse \(m\) est la lecture de l'espèce matérielle dans la carte choisie. Ces valeurs sont des sorties antérieures à cette réalisation. La matrice \(A_D\) est hermitienne et \(A_Dg^I\) est antihermitienne. Le terme cinétique symétrique possède donc la convention réelle correspondant à \eqref{vii:vii:eq:clifford-material-explicito}. Dans cette signature, l'équation sans torsion a pour terme principal \(\hbar g^I D_I\) ; la notation équivalente \(i\hbar\Gamma^I D_I\), avec \(\Gamma^I=-i g^I\), emploie la relation de Clifford de signature opposée. Cette précision fixe conjointement l'adjoint, la signature et le facteur imaginaire de l'action.

Les champs matériels et la cotétrade restent fixes dans la variation
de la connexion. Une préparation issue de l'extension minimale de
la section~\ref{vii:vii:sec:corriente-extension-minima} peut déterminer
la donnée initiale du champ ou de l'état. Chaque action conserve sa
variation propre : si l'on substitue dans une action un champ minimisant
dépendant de la connexion, on applique en outre sa différentielle composée,
ou le théorème de l'enveloppe sous ses hypothèses de stationnarité.
L'action~\eqref{vii:vii:eq:accion-dirac-material} réalise ici un
couplage affine en la contorsion.

\begin{proposition}[Courant spinoriel à normalisation variationnelle]
\label{vii:vii:prop:corriente-dirac-explicita}
Soit
\(B^{IJK}=\bar\psi g^{[I}g^Jg^{K]}\psi\), où les crochets
désignent l'antisymétrisation normalisée. Le courant défini dans
\eqref{vii:vii:eq:corriente-variacional} par l'action
\eqref{vii:vii:eq:accion-dirac-material} est
\begin{equation}
 \sigma_{JK}
 =-\frac{\hbar}{2}\,
   \bar\psi\{g^I,\mathbb B_{JK}\}\psi\,\eta_I
 =-\frac{\hbar}{2}\,B^I{}_{JK}\eta_I.
 \label{vii:vii:eq:corriente-dirac-trivector}
\end{equation}
Il est indépendant de la contorsion lorsque \(e\) et \(\psi\) sont
fixés. L'inversion de la section~\ref{vii:vii:subsec:reconstruccion-torsion}
s'applique donc à ce courant.
\end{proposition}
\begin{proof}
La variation de \(\Omega\) est \(\delta\Omega=\frac12\delta\omega^{JK}\mathbb B_{JK}\). Les deux termes cinétiques donnent respectivement le produit avec \(g^I\) à gauche et à droite. En conséquence,
\[
 \delta_\omega S_D
 =\frac{\hbar}{4}\int_M\delta\omega^{JK}\wedge\eta_I\,
                 \bar\psi\{g^I,\mathbb B_{JK}\}\psi.
\]
La comparaison avec \eqref{vii:vii:eq:corriente-variacional} fixe le
facteur et le signe de \eqref{vii:vii:eq:corriente-dirac-trivector}.
En développant l'anticommutateur et en utilisant la relation de Clifford,
les contractions de degré un s'annulent, et il reste
\(\{g^I,\mathbb B_{JK}\}=g^{[I}g_Jg_{K]}\).
L'action est linéaire en \(\Omega\), ce qui prouve la dernière
affirmation.
\end{proof}

\subsection{Contraction effective et coefficient du second moment}
\label{vii:vii:subsec:contraccion-espinorial}

La forme quadratique lorentzienne du trivecteur s'écrit
\begin{equation}
 q(B)=\frac1{3!}B^{IJK}B_{IJK}
 =-(B^{012})^2-(B^{013})^2-(B^{023})^2+(B^{123})^2.
 \label{vii:vii:eq:cuadratica-trivector}
\end{equation}
Cette contraction conserve la signature du module spinoriel. Son second moment dépend de l'état et de la réalisation des produits de champs.

\begin{theorem}[Interaction effective du courant spinoriel]
\label{vii:vii:thm:coeficiente-torsional-espinorial}
Avec l'action matérielle~\eqref{vii:vii:eq:accion-dirac-material},
l'élimination algébrique de la contorsion donne
\begin{equation}
 \mathcal L_{\rm spin}^{\rm eff}
 =C_\gamma q(B)\operatorname{vol}_e,\qquad
 C_\gamma=\frac{3\kappa c\hbar^2}{16}
                      \frac{\gamma^2}{1+\gamma^2}.
 \label{vii:vii:eq:coeficiente-cuadratico-espinorial}
\end{equation}
Le coefficient s'obtient à partir du courant et de l'inverse de Holst
avec les normalisations de l'action précédente.
\end{theorem}
\begin{proof}
Par \eqref{vii:vii:eq:accion-efectiva-spin}, il suffit d'évaluer \(-\frac14K_*^{IJ}\wedge\sigma_{IJ}\), avec
\[
 Y^{IJ}=\kappa c\frac{\gamma^2}{1+\gamma^2}
           (-\star+\gamma^{-1}I)\sigma^{IJ},\qquad
 K_*=\mathsf C_e(Y).
\]
Ici, \(\mathsf C_e\) est la composition explicite de
\eqref{vii:vii:eq:torsion-explicita} avec
\eqref{vii:vii:eq:contorsion-componentes}. La linéarité de ces deux
inverses réduit le calcul aux parties \(-\star\) et \(I\).
Pour donner la contraction complète, nous ordonnons les composantes du
trivecteur sous la forme \((012,013,023,123)\), extrayons les facteurs
\(\kappa c\hbar^2\gamma^2/(1+\gamma^2)\) et substituons
\(\sigma_{JK}=-\frac12B^I{}_{JK}\eta_I\). Les deux matrices
de la forme quadratique restante sont
\begin{equation}
 \mathcal B_{-\star}
 =\frac3{16}\operatorname{diag}(-1,-1,-1,1),\qquad
 \mathcal B_I=0_{4\times4}.
 \label{vii:vii:eq:matrices-contraccion-espinorial}
\end{equation}
Ces matrices s'obtiennent en utilisant uniquement \(e^I\wedge\eta_J=\delta_J^I\operatorname{vol}_e\) : pour chaque composante \(B^{abc}\), on forme \(\sigma\), on calcule \(Y\), puis \(T^I=\sum_J\iota_{E_J}Y^{IJ}
 -\frac14e^I\wedge\sum_{L,J}\iota_{E_L}\iota_{E_J}Y^{LJ}\), et enfin \(K_{IJK}=(T_{KIJ}-T_{IJK}-T_{JKI})/2\). La contraction donne les quatre entrées diagonales présentées ; les six polarisations entre composantes distinctes sont nulles pour les deux opérateurs. C'est le calcul de tous les coefficients des deux formes quadratiques. La partie \(\gamma^{-1}I\) s'annule et la partie \(-\star\) produit exactement \eqref{vii:vii:eq:cuadratica-trivector}, avec le coefficient de \eqref{vii:vii:eq:coeficiente-cuadratico-espinorial}.
\end{proof}

\subsection{Réalisation fermionique finie et second moment de l'état}
\label{vii:vii:subsec:segundo-momento-car}

Nous fixons une cellule comobile périodique, un repère orthonormé adapté à son
observateur et le mode spatial homogène des quatre composantes
spinorielles. Cette réalisation finie a pour espace d'états
\(\mathcal F=\bigoplus_{n=0}^4\Lambda^n\mathbb C^4\).
Soient \(a_r,a_r^\dagger\), \(r=1,\ldots,4\), ses opérateurs
canoniques, avec
\(\{a_r,a_s^\dagger\}=\delta_{rs}\) et
\(\{a_r,a_s\}=0\). Le vide d'occupation fixe l'ordre
normal. Pour une matrice \(M\), nous écrivons
\[
 \mathrm d\Gamma(M)=\sum_{r,s}a_r^\dagger M_{rs}a_s,
 \qquad \widehat N=\mathrm d\Gamma(I_4).
\]
La lettre \(\Gamma\) de cette notation désigne la seconde
quantification ; l'holonomie nonadique conserve sa propre notation.

Les quatre matrices des bilinéaires sont calculées directement :
\begin{equation}
 \begin{array}{c|cccc}
 abc&012&013&023&123\\ \hline
 M_{abc}=A_Dg^ag^bg^c
       &iJ\otimes R&iJ\otimes Z&iI_2\otimes J&iZ\otimes J
 \end{array}
 \label{vii:vii:eq:matrices-bilineales-espin}
\end{equation}
Toutes sont hermitiennes, ont une trace nulle et satisfont \(M_{abc}^2=I_4\). L'observable quartique correspondant à \eqref{vii:vii:eq:cuadratica-trivector}, dans cet ordre normal déclaré, est
\begin{equation}
 \begin{split}
 \widehat{\mathscr Q}_{\rm sp}
 &=\sum_{a<b<c}\eta_{aa}\eta_{bb}\eta_{cc}
       :\!\mathrm d\Gamma(M_{abc})^2\!:\\
 &=\sum_{a<b<c}\eta_{aa}\eta_{bb}\eta_{cc}
       \mathrm d\Gamma(M_{abc})^2+2\widehat N.
 \end{split}
 \label{vii:vii:eq:observable-cuadratico-espin}
\end{equation}
L'ordre normal fait partie de la réalisation finie spécifiée ;
sa définition rend explicite la soustraction des contractions
d'une seule occupation.

\begin{proposition}[Évaluation par les relations canoniques]
\label{vii:vii:prop:operador-espin-car}
L'opérateur \(\widehat{\mathscr Q}_{\rm sp}\) est autoadjoint,
conserve chaque secteur d'occupation et a pour restrictions
\begin{equation}
 \widehat{\mathscr Q}_{\rm sp}|_{\Lambda^0\oplus\Lambda^1}=0,
 \qquad
 \widehat{\mathscr Q}_{\rm sp}|_{\Lambda^3}=4I_4,
 \qquad
 \widehat{\mathscr Q}_{\rm sp}|_{\Lambda^4}=8I_1.
 \label{vii:vii:eq:sectores-espin-car}
\end{equation}
Dans la base de \(\Lambda^2\mathbb C^4\) dont les sous-ensembles d'occupation sont \((12,13,23,14,24,34)\), sa matrice est
\begin{equation}
 \begin{pmatrix}
 0&0&0&0&0&-4\\
 0&2&0&0&6&0\\
 0&0&2&-6&0&0\\
 0&0&-6&2&0&0\\
 0&6&0&0&2&0\\
 -4&0&0&0&0&0
 \end{pmatrix}.
 \label{vii:vii:eq:sector-dos-espin-car}
\end{equation}
\end{proposition}
\begin{proof}
La relation \(a_s a_t^\dagger=\delta_{st}-a_t^\dagger a_s\)
appliquée à \(\mathrm d\Gamma(M)^2\) sépare sa contraction à
une occupation et donne
\[
 :\!\mathrm d\Gamma(M)^2\!:
 =\mathrm d\Gamma(M)^2-\mathrm d\Gamma(M^2).
\]
Comme les trois premiers signes de
\eqref{vii:vii:eq:matrices-bilineales-espin} sont négatifs et le dernier
positif, la somme des contractions est \(-2\widehat N\).
Cela prouve \eqref{vii:vii:eq:observable-cuadratico-espin}.
L'autoadjonction et la conservation de l'occupation découlent de
ces mêmes expressions.

Dans \(\Lambda^0\), tous les termes s'annulent. Dans \(\Lambda^1\), \(\mathrm d\Gamma(M)=M\), et les carrés ont pour somme \(-2I_4\), qui s'annule avec \(2\widehat N\). Dans \(\Lambda^4\), \(\mathrm d\Gamma(M)=\operatorname{tr}M=0\), ce qui laisse huit. L'identification par complément entre \(\Lambda^3\mathbb C^4\) et le dual de \(\mathbb C^4\) envoie \(\mathrm d\Gamma(M)\) sur \((\operatorname{tr}M)I-M^{\mathsf T}=-M^{\mathsf T}\). Leurs carrés ont à nouveau pour somme \(-2I_4\), tandis que \(2\widehat N=6I_4\) ; la restriction est donc \(4I_4\). Enfin, l'application des quatre matrices explicites aux six produits extérieurs de la base indiquée produit \eqref{vii:vii:eq:sector-dos-espin-car}. L'opérateur est ainsi évalué dans les seize dimensions de \(\mathcal F\).
\end{proof}

Pour un opérateur d'état \(\varrho\geq0\),
\(\operatorname{Tr}\varrho=1\), nous définissons
\(q_\varrho=\operatorname{Tr}(\varrho\widehat{\mathscr Q}_{\rm sp})\).
Sa dépendance à l'égard des moments d'ordre deux est explicite :
\begin{equation}
 q_\varrho=\sum_{a<b<c}\eta_{aa}\eta_{bb}\eta_{cc}
 \left\{
  \langle\widehat B_{abc}\rangle_\varrho^2
  +\operatorname{Var}_\varrho(\widehat B_{abc})
  -\langle\widehat N\rangle_\varrho
 \right\},\qquad
 \widehat B_{abc}=\mathrm d\Gamma(M_{abc}).
 \label{vii:vii:eq:covarianza-espin-material}
\end{equation}
La covariance et la soustraction normale sont conservées même lorsque
les courants moyens s'annulent. En particulier, si \(P_3\) est la
projection sur \(\Lambda^3\mathbb C^4\),
\begin{equation}
 \varrho_3=\tfrac14P_3,\qquad
 \langle\widehat B_{abc}\rangle_{\varrho_3}=0,\qquad
 q_{\varrho_3}=4>0.
 \label{vii:vii:eq:testigo-positivo-espin}
\end{equation}
C'est un état positif de trace un qui constitue un témoin explicite
du domaine positif de la fonctionnelle. Le secteur de deux occupations
contient également des valeurs négatives ; par exemple, le vecteur
\((|12\rangle+|34\rangle)/\sqrt2\) a pour valeur \(-4\).
Le signe de l'intensité procède ainsi de la contraction et de
l'état, et reste distinct de la sélection des conditions cosmologiques.

\subsection{Densité, pression et conservation de l'amplitude comobile}
\label{vii:vii:subsec:amplitud-espinorial-comovil}

Dans la réalisation spatialement homogène, nous fixons
\[
 e^0=N(t)c\,\mathrm dt,\qquad e^j=a(t)\,\mathrm dx^j,
 \qquad \mathcal V(t)=a(t)^3V_c,
\]
avec une fonction lapse \(N(t)>0\), un facteur d'échelle sans dimension \(a(t)>0\)
et un volume comobile \(V_c>0\). Le cadre et le mode périodique de la
sous-section~\ref{vii:vii:subsec:segundo-momento-car} restent déclarés.
La normalisation du champ homogène est
\(\widehat\psi=\mathcal V^{-1/2}(a_1,a_2,a_3,a_4)^{\mathsf T}\).
Chaque bilinéaire porte un facteur \(\mathcal V^{-1}\), et son produit
normalement ordonné porte \(\mathcal V^{-2}\). La forme temporelle
cinétique symétrique conserve cette normalisation canonique : les termes
provenant de la dérivation de \(\mathcal V^{-1/2}\) s'annulent entre
ses deux termes.

\begin{theorem}[Amplitude torsionnelle comme fonctionnelle de l'état]
\label{vii:vii:thm:amplitud-material-explicita}
En maintenant fixes l'état canonique et la réalisation du produit
pendant les variations métriques, le terme effectif de spin a pour
action réduite
\begin{equation}
 S_{\rm spin}^{\rm hom}
 =\int\!\mathrm dt\,N(t)
       \frac{cC_\gamma}{a(t)^3V_c}\,q_\varrho.
 \label{vii:vii:eq:accion-homogenea-espin}
\end{equation}
Sa densité d'énergie et sa pression isotrope sont
\begin{equation}
 \rho_{\rm spin}=p_{\rm spin}
 =-\frac{cC_\gamma}{V_c^2}\,q_\varrho a^{-6}.
 \label{vii:vii:eq:densidad-presion-espin}
\end{equation}
Dans toute évolution qui conserve \(q_\varrho\), l'amplitude constante correspondante est
\begin{equation}
 \boxed{\displaystyle
 s_0[\varrho,V_c]
 =\frac{3\kappa c^2\hbar^2}{16}
    \frac{\gamma^2}{1+\gamma^2}
    \frac{\operatorname{Tr}
       (\varrho\widehat{\mathscr Q}_{\rm sp})}{V_c^2}.}
 \label{vii:vii:eq:amplitud-s0-espinorial}
\end{equation}
\end{theorem}
\begin{proof}
L'intégration spatiale de
\eqref{vii:vii:eq:coeficiente-cuadratico-espinorial} multiplie le
facteur \(\mathcal V^{-2}\) par
\(Nc\mathcal V\,\mathrm dt\), ce qui donne
\eqref{vii:vii:eq:accion-homogenea-espin}. En variables canoniques,
la variation du lapse définit l'énergie par
\(\delta_N S_{\rm m}=-\int E\,\delta N\,\mathrm dt\).
Par conséquent,
\[
 E_{\rm spin}=-\frac{cC_\gamma q_\varrho}{a^3V_c},
 \qquad \rho_{\rm spin}=\frac{E_{\rm spin}}{a^3V_c}.
\]
La variation spatiale isotrope satisfait \(\delta_a S_{\rm m}=\int3Na^2V_c p\,\delta a\,\mathrm dt\). Appliquée à \eqref{vii:vii:eq:accion-homogenea-espin}, elle donne
\[
 3Na^2V_c p_{\rm spin}
 =-\frac{3NcC_\gamma q_\varrho}{a^4V_c}.
\]
Cela prouve simultanément le signe, la pression et la densité de \eqref{vii:vii:eq:densidad-presion-espin}. Puisque \(q_\varrho\) est conservé, le coefficient de \(-a^{-6}\) est constant et l'on obtient \eqref{vii:vii:eq:amplitud-s0-espinorial}. La dimensionnalité est cohérente : \([cC_\gamma]=\text{énergie}\,L^3\), de sorte que \(s_0\) a la dimension d'une densité d'énergie.
\end{proof}

\begin{corollary}[Domaine conservé d'amplitude positive]
\label{vii:vii:cor:dominio-positivo-conservado}
Soit une évolution unitaire de \(\mathcal F\) qui préserve
\(\Lambda^3\mathbb C^4\). Pour tout état initial de trace un
à support dans ce secteur, on a \(q_{\varrho(t)}=4\), et
\begin{equation}
 s_0=\frac{3\kappa c^2\hbar^2}{4V_c^2}
                   \frac{\gamma^2}{1+\gamma^2}>0.
 \label{vii:vii:eq:amplitud-positiva-sector-tres}
\end{equation}
\end{corollary}
\begin{proof}
L'invariance du secteur conserve le support de l'état ; dans celui-ci,
\eqref{vii:vii:eq:sectores-espin-car} donne
\(\widehat{\mathscr Q}_{\rm sp}=4I\). La conservation de la
trace fixe alors son espérance à quatre, indépendamment de
l'évolution interne de ses cohérences. La positivité des
couplages, \(V_c>0\) et \(\gamma\ne0\) prouvent la formule
et son signe.
\end{proof}

Pour des états généraux, la conservation de
\(\langle\widehat N\rangle\) détermine la dilution de la densité
d'occupation, \(n(t)=\langle\widehat N\rangle/(a^3V_c)\).
La conservation du moment quartique est vérifiée séparément.
Si \(i\hbar\dot\varrho=[H(t),\varrho]\), alors
\begin{equation}
 \frac{\mathrm d q_\varrho}{\mathrm dt}
 =\frac{i}{\hbar}\operatorname{Tr}
       \bigl(\varrho[H(t),\widehat{\mathscr Q}_{\rm sp}]\bigr).
 \label{vii:vii:eq:conservacion-momento-cuartico}
\end{equation}
La commutation, l'invariance d'un secteur où l'observable est scalaire, ou une conservation démontrée pour l'état concret fournissent les conditions pertinentes. Si le moment évolue, \eqref{vii:vii:eq:densidad-presion-espin} conserve son évaluation instantanée ; son bilan inclut
\[
 \dot\rho_{\rm spin}+6H\rho_{\rm spin}
 =-\frac{cC_\gamma}{V_c^2}\,a^{-6}\dot q_\varrho,
 \qquad H=\dot a/a
\]
en temps propre. Ce terme enregistre l'échange avec les
autres degrés de liberté matériels.

La composition est déterminée : représentation spinorielle du
transport, action matérielle, courant variationnel, inversion de
Cartan--Holst, second moment et densité normalisée. Le domaine
\(s_0>0\) de la section~\ref{vii:vii:sec:dilucion-rebote} dispose ainsi
d'une fonctionnelle évaluable et d'états témoins explicites. Son utilisation dans
la dynamique cosmologique conserve en outre les conditions matérielles,
géométriques et de conservation énoncées dans cette section.
